\chapter{Conclusion and Future Work}
\label{ch:conclusion}

\section{Summary of Contributions}
\label{sec:conclusion_summary}

This thesis presented a privacy-preserving, federated anomaly detection
framework for distributed microservice architectures. The central novelty is
the end-to-end integration of edge-local LSTM-AE detection, FedAvg, DP-SGD,
communication compression, causal RCA, and adaptive thresholding into one
evaluated monitoring pipeline. The contribution is not an isolated model
component, but a complete privacy-aware monitoring architecture with measured
accuracy, communication, privacy, RCA, and scalability behavior.

The thesis delivers six concrete outcomes: (1)~a federated LSTM-AE whose
detection performance is reported in~\S\ref{sec:eval_detection};
(2)~differential privacy with empirically characterized gradient clipping and
noise trade-offs (\S\ref{subsec:eval_dp_impact}); (3)~causal-graph-based root
cause analysis (\S\ref{sec:eval_rca}); (4)~an exploratory Q-learning
threshold study showing partial stabilization and cold-start instability
(\S\ref{sec:eval_threshold}); (5)~an adaptive compression pipeline
(\S\ref{sec:eval_fl_comm}); and (6)~fault-tolerant edge design with two node
profiles (\S\ref{sec:reliability_overview}).

\section{Conclusion}
\label{sec:conclusion_conclusion}

The evaluation confirms the main thesis claim: privacy-aware federated anomaly
detection can be practical for edge-monitored microservice systems when model
training, communication, and RCA are designed together. Detection results
(\S\ref{sec:eval_detection}) show that the federated LSTM-AE outperforms the
centralized baseline on SMD, with the performance gap analyzed in
\S\ref{sec:eval_discussion}. Communication remains modest: binary serialization
reduces payload size substantially (\S\ref{sec:eval_fl_comm}), and FedAvg
aggregation scales approximately linearly to 1\,000 simulated nodes
(\S\ref{sec:eval_scalability}).

The privacy and adaptivity results refine the claim rather than simply
confirming it. DP-SGD remains near baseline for mild noise, but F1 drops
sharply at $\sigma{=}0.001$ with default clipping; reducing $C$ to 0.01
recovers utility, exposing a clipping--privacy tension discussed in
\S\ref{subsec:eval_clipping}. The RCA engine provides operationally useful
Top-3 rankings on small topologies but degrades on deeper dependency graphs
(\S\ref{sec:eval_rca}). Q-learning thresholding reduces final alert volume
after extended interaction but suffers from cold-start instability
(\S\ref{sec:eval_threshold}). The framework is therefore strongest as a
detection, privacy, compression, and RCA pipeline, with adaptive thresholding
left as a promising but not yet deployment-ready research path.

\section{Limitations}
\label{sec:conclusion_limitations}

Several limitations remain. A 3-seed retraining
(\S\ref{sec:eval_detection}) shows that AUC-ROC is stable
($0.932 \pm 0.017$) but F1 varies with threshold calibration
($0.743 \pm 0.077$); this variance is not yet quantified for the DP
or RCA benchmarks.
The tabular Q-learning agent eventually stabilizes, but its early alert burst
and negative cumulative reward show that warm starts and improved reward
shaping are still required. FedAvg assumes approximately IID local datasets
\cite{li2020federated}, so highly heterogeneous production workloads may
require FedProx or scaffold methods. RCA remains weaker on deep dependency
chains, where Top-1 accuracy matches random selection on the 50-node Alibaba
topology. Finally, tight clipping recovers detection utility but reduces the
absolute perturbation $\sigma C$, so formal DP accounting should be
complemented with empirical reconstruction and membership-inference testing.

\section{Future Work}
\label{sec:conclusion_future}

Future work should extend the framework along four prioritized directions:

\begin{enumerate}
    \item \textbf{Empirical privacy validation} (highest priority): run
    membership inference attacks~\cite{shokri2017membership} against trained
    LSTM-AE models under different $(\sigma, C)$ configurations to validate
    whether the theoretical $\varepsilon$ bounds translate to practical attack
    resistance.

    \item \textbf{DQN-based threshold control}: replace tabular Q-learning
    with a Deep Q-Network using a 4-dimensional continuous state space
    $[\tau, \text{FPR}, \text{FNR}, \text{SLO}]$ and an experience replay
    buffer of size 10\,000; expected to eliminate the cold-start alert burst
    observed at minute~360 by enabling warm-start learning from stored
    incidents.

    \item \textbf{Hierarchical federated learning}: introduce a two-tier
    coordinator hierarchy (cluster-level sub-coordinators + global
    coordinator) to reduce per-round communication by approximately
    $K\times$, where $K$ is cluster size; validate on the Alibaba trace with
    500+ nodes.

    \item \textbf{Richer RCA graphs}: combine metrics, logs, traces, and
    deployment metadata into a unified causal graph; evaluate whether
    multi-signal evidence improves Top-1 accuracy on deep topologies beyond
    the current 28\% on the 50-node Alibaba graph.
\end{enumerate}

In summary, the thesis demonstrates that privacy-aware federated anomaly
detection is feasible for edge-monitored microservice systems, while also
showing that deployment quality depends on the interfaces between detection,
privacy, communication, RCA, and threshold control. The next technical step is
a stronger validation loop: extending seed variance to DP and RCA benchmarks,
attack-based privacy testing, warm-started threshold policies, and larger
labeled RCA traces.
